
\documentclass{article}[12pt]

\usepackage[utf8]{inputenc}
\usepackage{amsmath}
\usepackage{amsfonts}
\usepackage{amssymb}
\usepackage{amsthm}
\usepackage{hyperref}
\usepackage{tikz}
\usepackage{graphicx}
\usepackage{cleveref}
\usepackage{natbib}
\usepackage{enumitem}
\usepackage{booktabs}
\usepackage{multirow}
\usepackage{multicol}
\usepackage{array}
\usepackage{float}
\usepackage{caption}

\captionsetup[table]{width=0.95\textwidth, font=small}
\captionsetup[figure]{width=0.95\textwidth, font=small}

\hypersetup{
	colorlinks=true,
	linkcolor=black,
	citecolor=black}

%%% TIKZ graphics
\usetikzlibrary{positioning, calc, shapes.geometric, shapes.multipart, shapes, arrows.meta, arrows, 	decorations.markings, external, trees}

\tikzstyle{Arrow} = [
thick,
decoration={
	markings,
	mark=at position 0.99 with {  % fix bug with tikz arrow head interference with decoration
		\arrow[thick]{latex}
	}
},
shorten >= 3pt, preaction = {decorate}
]
\tikzstyle{BiArrow} = [
thick,
decoration={
	markings,
	mark=at position 0.99 with {  % fix bug with tikz arrow head interference with decoration
		\arrow[thick]{latex}
	}
},
shorten >= 3pt, preaction = {decorate}
]
\tikzstyle{RedArrow} = [
red, 
line width = 0.5mm,
decoration={
	markings,
	mark=at position 0.99 with {
		\arrow[line width = 0.5mm,red]{latex}
	}
},
shorten >= 3pt, preaction = {decorate}
]

\tikzstyle{BlueArrow} = [
blue, 
line width = 0.5mm,
decoration={
	markings,
	mark=at position 0.99 with {
		\arrow[line width = 0.5mm,blue]{latex}
	}
},
shorten >= 3pt, preaction = {decorate}
]

\tikzstyle{Dashed} = [
thick,dashed,
decoration={
	markings,
	mark=at position 0.99 with {  % fix bug with tikz arrow head interference with decoration
		\arrow[thick]{latex}
	}
},
shorten >= 3pt, preaction = {decorate}
]

\tikzset{
    BiArrow/.style={
        thick,
        decoration={
            markings,
            mark=at position 0.99 with {
                \arrow[thick]{latex}
            },
            mark=at position 0 with {
                \arrow[thick]{latex}
            }
        },
        shorten >= 3pt,
        preaction = {decorate}
    }
}



\newtheorem{remark}{Remark}
\AfterEndEnvironment{remark}{\noindent\ignorespaces}
\newtheorem{assumption}{Assumption}
\AfterEndEnvironment{assumption}{\noindent\ignorespaces}
\newtheorem{proposition}{Proposition}
\AfterEndEnvironment{proposition}{\noindent\ignorespaces}
\newtheorem{theorem}{Theorem}
\AfterEndEnvironment{theorem}{\noindent\ignorespaces}
\newtheorem{corollary}{Corollary}
\AfterEndEnvironment{corollary}{\noindent\ignorespaces}
\newtheorem{condition}{Condition}
\AfterEndEnvironment{condition}{\noindent\ignorespaces}
\newtheorem{lemma}{Lemma}
\AfterEndEnvironment{lemma}{\noindent\ignorespaces}

% \AtAppendix{\counterwithin{lemma}{section}}

\theoremstyle{definition}
\newtheorem{exmp}{Example}
\AfterEndEnvironment{exmp}{\noindent\ignorespaces}
\newtheorem{design}{Design}
\AfterEndEnvironment{design}{\noindent\ignorespaces}
\newtheorem{defi}{Definition}
\AfterEndEnvironment{defi}{\noindent\ignorespaces}

\def \bc {\boldsymbol{c}}
\def \bz {\boldsymbol{z}}
\def \bZ {\boldsymbol{Z}}
\def \bA {\boldsymbol{A}}
\def \ba {\boldsymbol{a}}
\def \bX {\boldsymbol{X}}
\def \bx {\boldsymbol{x}}
\def \bY {\boldsymbol{Y}}
\def \wY {\widetilde{Y}}
\def \bI {\boldsymbol{I}}
\def \bS {\boldsymbol{S}}
\def \bR {\boldsymbol{R}}
\def \bO {\boldsymbol{O}}
\def \bG {\boldsymbol{G}}
\def \bE {\boldsymbol{E}}
\def \bV {\boldsymbol{V}}
\def \bf {\boldsymbol{f}}
\def \bD {\boldsymbol{D}}
\def \cR {\mathcal{R}}

\def \EX {\mathbb{E}}

\def \beps {\boldsymbol{\varepsilon}}

\newcommand{\indep}{\mathrel{\perp\!\!\!\perp}}
\newcommand{\cindep}[3]{#1 \indep #2 \mid #3}

\addtolength{\oddsidemargin}{-.5in}%
\addtolength{\evensidemargin}{-1in}%
\addtolength{\textwidth}{1in}%
\addtolength{\textheight}{1.7in}%
\addtolength{\topmargin}{-1in}%


\usepackage[status=draft,inline,index]{fixme}
%\usepackage[status=final,inline,index]{fixme}
\fxsetup{theme=color,margin=false,mode=multiuser}
\FXRegisterAuthor{bw}{bwe}{BW}


\title{Sensitivity analysis for contamination in cluster randomized trials with interference}

\author{Bar Weinstein}
% date is current date
\date{\today}

\begin{document}
\maketitle


\section{Porblem setup}
Consider a CRT with one stage randomization. Assume there are clusters
$C_1, \ldots, C_K$ in the population, where $K$ is the total number of clusters.
Let $n_k$ be the size of cluster $C_k$, and $N = \sum_{k=1}^K n_k$ be the total population size.
Assume researchers randomize binary treatments at the cluster level. 
Let $Z_k$ be the treatment assignment of cluster $C_k$, where $Z_k=1$ if cluster $C_k$ is treated, and $0$ otherwise.
Assume that $P(Z_k = 1) = p_z \in (0,1)$ for all $k=1,\ldots,K$.

The observed network $\widetilde{\bA}$ consists of $K$ disjoint clusters, each 
assumed to be fully connected. The true network $\bA$ may contain edges between clusters,
leading to contamination across clusters.

Let $Z_{ik}$ be the treatment assignment of unit $i$ in cluster $C_k$.
Assume the binary exposure mapping
\begin{assumption}[Exposure mapping]
\label{ass:expos_map}
    Let $F(\bZ_{-i},\bA) =\mathbb{I}\left\{\sum_{j\neq i} Z_jA_{ij} > 0\right\}$. For any $\bz, \bz' \in \mathcal{Z}$, if $z_i = z'_i$ and $F(\bz_{-i},\bA) = F(\bz'_{-i},\bA)$, then $Y_i(\bz) = Y_i(\bz')$, 
\end{assumption}
where $\mathcal{Z} \subseteq \{0,1\}^N$ is the space of all possible treatments allocations with 
non-zero probability.
onsequently, each unit has four effective potential outcomes,
$Y_i(z,f)$ for $z,f\in\{0,1\}$, indicating whether unit $i$ is treated and whether at least one of its neighbors is treated.
The observed exposures are $\widetilde{F}_i = F(\bZ_{-i}, \widetilde{A})$.
By the CRT design, $\widetilde{F}_i = Z_i$ with probability one, 
as units within a cluster share the same treatment assignment, and the observed network 
consists of disjoint clusters.

Assume that the observed outcomes $Y_i$ are generated from one of the potential outcomes by:
\begin{assumption}[Consistency]
\label{ass:consis}
$Y_i = \sum_{z,f \in \{0,1\}}Y_i(z,f) \mathbb{I}\left\{Z_i=z, F_i=f\right\}$ .    
\end{assumption}

The causal estimand is the sample average treatment effects. Using our setup and notations it reduced to
\begin{equation}
    \label{eq:ate}
    ATE = \frac{1}{N}\sum_{i=1}^N Y_i(1, 1) - Y_i(0, 0) = 
    \frac{1}{N}\sum_{k=1}^{K}\sum_{i \in C_k} Y_i(1, 1) - Y_i(0, 0).
\end{equation}
The most common estimator for the ATE in CRTs is the Horvitz-Thompson (HT) estimator,
\begin{equation}
    \label{eq:ate_ht}
    \begin{aligned}
        \widehat{ATE} 
        &= 
        \frac{1}{N}
        \sum_{k=1}^{K}
        \sum_{i \in C_k}
        \left[
        \frac{
        \mathbb{I}\{Z_k=1\}Y_i}{p_z} 
        -
        \frac{
        \mathbb{I}\{Z_k=0\}Y_i}{1 - p_z} 
        \right]
        \\ &=
        \frac{1}{N}
        \sum_{i=1}^{N}
        \left[
        \frac{  
        \mathbb{I}\{Z_i=1\}Y_i}{p_z}
        -
        \frac{
        \mathbb{I}\{Z_i=0\}Y_i}{1 - p_z}
        \right].
    \end{aligned}
\end{equation}
For units in a treated cluster, the observed exposure $\widetilde{F}_i=1$ agrees with the true exposure $F_i=1$.
Therefore
\begin{equation*}
    \EX_{\bZ} \left[\frac{
        \mathbb{I}\{Z_i=1\}Y_i}{p_z}
    \right]
    = Y_i(1,1).
\end{equation*}
However, for units in a control cluster, the observed exposure $\widetilde{F}_i=0$ 
may not agree with the true exposure $F_i$, which can be either $0$ or $1$ due to contamination.
Notice that $\mathbb{I}\{Z_i=0\} = \mathbb{I}\{\widetilde{F}_i=0\}$
with probability one as $\widetilde{F}_i=1$ only if $Z_i=1$.
Thus, 
\begin{equation*}
    \begin{aligned}
        \EX_{\bZ} \left[\frac{
            \mathbb{I}\{Z_i=0\}Y_i}{1 - p_z}
        \right]
        &= 
        \EX_{\bZ} \left[\frac{
            \mathbb{I}\{\widetilde{F}_i=0\}Y_i}{1 - p_z}
        \right]
        \\ &=
        \EX_{\bZ} \left[\frac{
            \mathbb{I}\{\widetilde{F}_i=0, F_i=0\}Y_i(0,0)}{1 - p_z}
        \right]
        \\&\quad +
        \EX_{\bZ} \left[\frac{
            \mathbb{I}\{\widetilde{F}_i=0, F_i=1\}Y_i(0,1)}{1 - p_z}
        \right]
        \\ &=
        Y_i(0,0) \Pr(F_i = 0 \mid \widetilde{F}_i=0)
        \\&\quad +
        Y_i(0,1) \Pr(F_i = 1 \mid \widetilde{F}_i=0).
    \end{aligned}
\end{equation*}
By Bayes' rule,
\begin{equation*}
    \begin{aligned}
        \Pr(F_i = 0 \mid \widetilde{F}_i=0) 
        &=
        \frac{
            \Pr(\widetilde{F}_i=0 \mid F_i=0) \Pr(F_i=0)
        }{
        \Pr(\widetilde{F}_i=0)
        }
        \\ &=
        \frac{
            1-\pi_i
        }{
          1-p_z
        }.
    \end{aligned}
\end{equation*}
Where we used the fact that $ \Pr(\widetilde{F}_i=0 \mid F_i=0) = 1$
as $\sum_{j\neq i}\widetilde{A}_{ij}Z_j \leq \sum_{j\neq i}A_{ij}Z_j$ with 
probability one. Therefore, $\widetilde{F}_i \leq F_i$ as well. Since 
the exposures are binary, $F_i=0$ implies $\widetilde{F}_i=0$ with probability one,
Also, $\Pr(\widetilde{F}_i=0) = \Pr(Z_i=0) = 1-p_z$ by the CRT design.
And we use the notation $\pi_i = \Pr(F_i=1)$.
Consequently,
\begin{equation*}
     \EX_{\bZ} \left[\frac{
            \mathbb{I}\{Z_i=0\}Y_i}{1 - p_z}
        \right] 
        = 
        \frac{
            (1- \pi_i)Y_i(0,0) + (\pi_i-p_z)Y_i(0,1)
        }{
            1-p_z
        }.
\end{equation*}
That enables us to express the expectation of the HT estimator as
\begin{equation*}
    \begin{aligned}
        \EX_{\bZ}[\widehat{ATE}] 
        &=
        \frac{1}{N}
        \sum_{i=1}^{N}
        \left[
        Y_i(1,1)
        -
        \frac{
            (1- \pi_i)Y_i(0,0) + (\pi_i-p_z)Y_i(0,1)
        }{
            1-p_z
        }
        \right]
        \\ &=
        ATE
        -
        \frac{1}{N}
        \sum_{i=1}^{N}
        \frac{
            \pi_i-p_z 
        }{
            1-p_z
        }
        \left(Y_i(0,1) - Y_i(0,0)\right).
    \end{aligned}
\end{equation*}
Notice that 
\begin{equation*}
    \begin{aligned}
        \pi_i &=
        \Pr\left(F_i=1\right)
        \\ &= 
        \Pr\left(F_i = 1 \mid Z_i=1\right) \Pr(Z_i=1)
        + \Pr\left(F_i = 1 \mid Z_i=0\right) \Pr(Z_i=0)
        \\ &=
        p_z + (1-p_z) \Pr\left(F_i = 1 \mid Z_i=0\right)
        \\ & \geq p_z.
    \end{aligned}
\end{equation*}
Therefore, if $Y_i(0,1) \geq Y_i(0,0)$ for all $i$,
then $\EX_{\bZ}[\widehat{ATE}] \leq ATE$, and if 
$Y_i(0,1) \leq Y_i(0,0)$ for all $i$,
then $\EX_{\bZ}[\widehat{ATE}] \geq ATE$.

\paragraph{Bias-corrected estimator.}
An unbiased estimator for the ATE can be obtained by correcting the bias of the HT estimator.
Define the bias-corrected estimator as  
\begin{equation}
    \label{eq:ate_corrected}
    \begin{aligned}
        \widehat{ATE}_{corrected} 
        &=
        \widehat{ATE} 
        +
        \frac{1}{N}
        \sum_{i=1}^{N}
        \frac{
            \pi_i - p_z
        }{
            1-p_z
        }
        \left[
            \frac{\mathbb{I}\{Z_i=0, F_i=0\}Y_i}{
                \Pr(Z_i=0, F_i=0)
            }
            -
            \frac{\mathbb{I}\{Z_i=0, F_i=1\}Y_i}{
                \Pr(Z_i=0, F_i=1)
            }
            \right].
    \end{aligned}
\end{equation}
The problem is that we don't have any information in the observed data about
$Y_i(0,1)$, as we don't know which units in the control clusters have $F_i=1$.
Assuming (conditional) edge independence in the population network,
we can express $\pi_i$ as a function of the individual edge probabilities. 
Using this specification, we can simulate contamination,
compute the probabilities $\Pr(Z_i=0, F_i=f)$ for $f\in\{0,1\}$,
obtain the units with $F_i=1$ in control clusters,
 and compute the bias-corrected estimator \eqref{eq:ate_corrected}.
This enables us to combine the simple (yet general) simulation-based sensitivity analysis approach 
with analytical expression of the bias-corrected estimator.



\end{document}
